\section{Proof of Theorem \ref{th:stockmeyer}: A Sub-Exponential Time Approximate Counter}
\label{sec:stockmeyer}

%\begin{theorem}
 %   For any constants $r$ and $\rho$, there is a randomized $\tilde{O}(\sqrt{C(P)})$-time $(r, \rho)$-approximate counter.
%\end{theorem}

To describe our algorithm, we assume that we are given a concurrent program $P$, 
which defines the tree $\mathcal{D}(P)$. 
As mentioned before, if $d$ is the size of the program $P$, 
this tree has height $d^2$ and an inspection of the construction shows that the degree of each node is at most $d^2$. 
However, every node need not have degree $d^2$. In order to make this tree uniform, we add dummy nodes to each node of $\mathcal{D}(P)$ so that the degree of each node becomes $d^2$. 
Call the resulting tree $\mathcal{H}(P)$. Note that $\mathcal{H}(P)$ 
is the complete $d^2$-ary tree of height $d^2$. 
(We will never explicitly construct $\mathcal{D}(P)$ nor $\mathcal{H}(P)$ during our algorithm; 
these are just defined implicitly). 



We define an \emph{inorder numbering} of $\mathcal{H}(P)$, which associates a subset of numbers to each node of $\mathcal{H}(P)$ recursively in the following way: We start at the root with a variable called count initialized to 0. For each $i$ from 1 to $d^2$, we recursively compute the inorder numbering of the $i^{th}$ subtree of the root. At the end of this computation, let the value of count be $n_i$. We increment $n_i$ by 1,
add $n_i+1$ to the inorder numbering of the root and move to the next iteration of the loop (i.e., to $i+1$). Note that the base case of the recursion is when we arrive at a leaf, in which case we will simply increment count by 1 and add this incremented value to the inorder numbering of the leaf and go back to its parent. In this way, we associate a subset of numbers to each node of $\mathcal{H}(P)$. 

For any node $u$, let $M(u)$ be the subset of numbers associated to $u$ by the inorder numbering. Let $S(u)$ and $B(u)$ be the smallest and biggest numbers appearing in $M(u)$ respectively. It is easy to see that $M(u)$ is a single number for each leaf
and otherwise, consists of exactly $d^2$ numbers. Furthermore, $M(u) \cap M(v) = \emptyset$
for any two distinct nodes $u,v$. 
Hence, from the construction, it follows that each number between $1$ and $N = (d^2)^{d^2} + (d^2)[{(d^2)^{d^2+1}} - 1]$ is assigned to exactly one node.

Now, for any pair of numbers $1 \le a \le b \le N$, let $[a,b]$ denote the interval of natural numbers between $a$ and $b$. Notice that if $M(u) \cap [a,b] \neq \emptyset$
and $M(v) \cap [a,b] \neq \emptyset$, then the least common ancestor $w$ of $u,v$ also satisfies $M(w) \cap [a,b] \neq \emptyset$. Hence, for each interval $[a,b]$, there
is a unique node $w_{a,b}$ such that $M(w_{[a,b]}) \cap [a,b] \neq \emptyset$ and 
every node $u$ satisfying $M(u) \cap [a,b] \neq \emptyset$ belongs to the subtree of $w_{[a,b]}$. 

Based on this notion of inorder numbering, we now present our estimation algorithm. Our estimation algorithm works on a parameter $\theta$ and consists of two acts: one where we operate under the assumption that the number of leaves of $\mathcal{D}(P) \le \theta$ and another where we operate under the assumption that the number of leaves of $\mathcal{D}(P) > \theta$.

\paragraph{Act One: Number of leaves of $\mathcal{D}(P) \le \theta$. }
In order to describe the first act of the algorithm, we need some notation and some small observations. Recall that any number between 1 and $N$ is assigned to exactly one node. Therefore, we can order the nodes of $\mathcal{H}(P)$ totally as follows:
$u \preceq v$ iff $S(u) < S(v)$. We now make the following series of observations.


\subparagraph*{Observation 1: } We can compute the inorder numbering $M(w)$ of any node as follows. We start by setting $u$ to be the root and a counter $x$ set to 0. Intuitively $x+1$ will denote the smallest number that can appear in the inorder numbering of any descendant of $u$.
At each step, we find the unique child $v$ of $u$ that contains $w$ as its descendant.
Suppose $v = w$. If $v$ is a leaf, then we simply return $\{x+1\}$.
If $v$ is an internal node at level $d-j$ (for some $j$), we return $\{x + N_{j-1} + 1, x + 2 N_{j-1} + 2, \dots, x + d^2 N_{j-1} + d^2\}$ where $N_{j-1} = (d^2)^{j-1} + d^2 [(d^2)^j - 1]$. Otherwise, if $v$ is the $i^{th}$ child for some $i \ge 1$ and $v$ is at level $d-j$ (for some $j$), then we set $x = x + (i-1) N_{j-1} + (i-1)$ and set $u$ to $v$ and continue. It can be easily checked that this procedure correctly computes $M(w)$ and 
runs in time $d^2$.

% Given any node $u$ of $\mathcal{H}(P)$ with inorder numbering $M(u)$, we can compute the inorder numbering $M(v)$ of any child $v$ of $u$ as follows. Indeed, let $v$ be the $i^{th}$ child of $u$ and let $x$ be the $(i-1)^{th}$ inorder number in $M(u)$ if $i-1 \ge 1$ and let $x$ be 0 otherwise. If $v$ is a leaf, then $M(v)$ is simply $\{x+1\}$. Otherwise, suppose $v$ is an internal node at level $d-j$.
% Let $x$ be the $(i-1)^{th}$ inorder number in $M(u)$ if $i - 1 \ge 1$ and let $x$ be 0 otherwise. Then, $M(v)$ is simply $\{x + N_{j-1} + 1, x + 2 N_{j-1} + 2, \dots, x + d^2 N_{j-1} + d^2\}$ where $N_{j-1} = (d^2)^{j-1} + d^2 (d^j) - d^2$.

\subparagraph*{Observation 2: } Given any node $u$ of $\mathcal{H}(P)$, we can check if 
there is a leaf of $\mathcal{D}(P)$ at the subtree rooted at $u$. Indeed, if $u$
corresponds to a dummy node, i.e., not a node of $\mathcal{D}(P)$, then
no such node will exist. On the other hand, if $u$ is a node of $\mathcal{D}(P)$
then indeed there will be a leaf of $\mathcal{D}(P)$ at the subtree rooted at $u$.
Hence, this check can be simply reduced to checking if $u$ is a node of $\mathcal{D}(P)$,
which can be done quite efficiently.

\subparagraph*{Observation 3: } Given any interval $I = [a,b]$, we can find the node 
$w_{I}$ as follows: We initially begin with the node $u$ set to the root and an interval $I'$ set to $[1,N]$. In each step, we will maintain the invariant that $w_I$ will be a descendant of $u$, $u = w_{I'}$ and also that $I \subseteq I'$. Indeed, this is true at the beginning.
Suppose at some step, we have reached $u$ with interval $I'$ such that $u = w_{I'}$
and $I \subseteq I'$. If $u$ is a leaf, then $w_I$ must be $u$ and so we are done.
Otherwise let $M(u)$ be $\{x_1,\dots,x_{d^2}\}$ and let $I' = [a',b']$. If $I$ is completely contained
inside any one of $I'_1 = [a'+1,x_1-1], I'_2 = [x_1+1,x_2-1], [x_2+1,x_3-1],\dots, I'_{d^2} = [x_{d^2-1}+1,b'-1]$, then we pick the unique $j$ such that $I$ is completely contained inside $I'_j$. Then we update $u$ to be the $j^{th}$ child of $u$ and update $I'$ to be $I'_j$. (Notice that the invariant is satisfied in this case). Otherwise,
it must be the case that $u = w_I$ and so we just output $u$. In this way, by going down the tree, in $d^2$ steps, we can find the node $w_I$.


Now, let $v_1,v_2,\dots$ be the leaves of $\mathcal{D}(P)$ (the \trust tree, not the complete $d^2$-ary tree) ordered according to the $\preceq$ ordering. Using the above mentioned observations and notations, we will now describe an algorithm which will find the nodes $v_1,v_2,\dots,v_{\theta}$. It does so in $\theta$ many phases, where in the $i^{th}$ phase it will find the node $v_i$.

The $i^{th}$ phase will operate in $d^2$ many subphases. At the start of each subphase
we will maintain a node $z$ of $\mathcal{D}(P)$ and an interval $I = [a,b]$ such that
$z = w_I$. (For the first subphase of the first phase, $z$ is taken to be the root and $I$ is taken to be $[1,N]$). For each subphase, we will always maintain the invariant that the node $v_i$ we are looking for will be a 
descendant of $z$, $S(v_i) \in I$ and $v_{i-1}$ will not be a descendant of $z$. (Notice that this invariant is definitely satisfied at the beginning of the first subphase of the first phase).

In the $j^{th}$ subphase, given the node $z$ (with corresponding inorder numbering $M(z)$) and the interval $I = [a,b]$, we proceed as follows. If $a = b$, then by our invariants, it follows that $v_i = z$ and so we return $z$. Otherwise, we loop over all of the $d^2$ children of $z$ (say $u_1,u_2,\dots,u_{d^2}$ such that $u_1 \preceq u_2 \preceq \dots \preceq u_{d^2}$). 
When looping over $u_k$, we first compute $M(u_k)$ using Observation 1. If $M(u_k) \cap I = \emptyset$, then we move to $u_{k+1}$ in the loop. Further, using Observation 2, we check if there is some node of $\mathcal{D}(P)$ that is present in the subtree rooted at $u_k$. If not, then also we move to $u_{k+1}$ in the loop. 
Else, it immediately follows that $v_i$ must be a descendant of $u_k$. 
Hence, we move to the $(j+1)^{th}$ subphase with the node $u_k$ and compute the new interval $I'$ as follows: Let $x$ be the $(k-1)^{th}$ number appearing in $M(z)$ if $k - 1 \ge 1$, otherwise let $x$ be $a$. Let $y$ be the $(k+1)^{th}$ number appearing in $M(z)$ if $k+1 \le d^2$, otherwise let $y$ be $b$. We then set $I'$ to be $[x,y]$. It can be easily verified that $u_k = w_{I'}$. Furthermore because $v_i$ must be a descendant of $u_k$, it follows that $S(v_i) \in I'$. Hence the invariants for the next subphase are satisfied.

Note that in each subphase, we are going one level down the tree. Hence there can be at most $d^2$ subphases in each phase. Furthermore, due to our invariants, it follows
that when the $i^{th}$ phase ends, we have arrived at the node $v_i$ with an interval $I$ of the form $[a,a]$. We then start the $(i+1)^{th}$ phase by setting our interval to be $[a+1,N]$ and computing the corresponding node $w_{a+1,N}$ using Observation 3. 
It easily follows that our invariant is satisfied at the beginning of the $(i+1)^{th}$ phase. This proves that the algorithm finds the nodes $v_1,v_2,\dots,v_{\theta}$.

Let us analyse the running time of this algorithm. Notice that each subphase of each phase
takes at most $O(d^4)$ time, in order to go over all of the possible children of $z$ and compute each of their inorder numberings. Since there are $d^2$ subphases in total and $\theta$ phases overall, it follows that the total amount of time spent in this act is $O(d^6 \theta)$. Note that if the number of leaves of $\mathcal{D}(P)$ is indeed
at most $\theta$, then our algorithm will stop with this act and output all of them.
Since it uses no randomization, in this case, it will output the exact number of leaves
of $\mathcal{D}(P)$ with zero variance.

\paragraph{Second Act: Number of leaves of $\mathcal{D}(P) > \theta$. }

Suppose we executed the first act of our algorithm to completion. It follows that there
are at least $\theta$ many leaves of $\mathcal{D}(P)$. At this point, we initiate the second act of our algorithm. Let $M = ((d^2)^{d^2+1})-1$ be the total number of nodes
of $\mathcal{H}(P)$. We take $z$ independent random walks and count the number of times we hit a maximal execution graph (say $y$). Then we output $(M/z) \cdot y$. 

Let $X$ be the random variable corresponding to the value that we output.
It is easy to see that the expected value of $X$ is the precise number
of leaves of $\mathcal{D}(P)$. Let us now count the variance of this approach, noting that
we are guaranteed that the number of leaves of $\mathcal{D}(P)$ is at least $\theta$
and therefore the probability of hitting a maximal execution graph along a random walk
is at least $p = \theta/M$. 

Let $Y$ be the random variable which counts the number of successful samples from $z$ independent random walks. Since this is essentially a binomial distribution it follows that
the variance of $Y$ is $zp(1-p)$. Hence, the variance of $X$ is the variance of
$\frac{M}{z} \cdot Y$ which is $\frac{M^2}{z^2} \cdot z p (1-p) = \frac{M^2 p (1-p)}{z}$. 
Let $\sigma$ be the square root of the variance of $X$. 

Now, suppose we are given numbers $r, \rho$ and we have to choose a $z$ so that the random variable $X$ is at most a factor of $r$ away from the total number of leaves of $\mathcal{D}(P)$ with probability at least $1-\rho$, i.e., we want to choose a $z$ so that $\mathbb{P}(X \ge r \mu )$ is at most $\rho$. Taking $z = \frac{1}{\rho} \cdot \frac{1}{(r-1)^2} \cdot \left(\frac{M}{\theta} - 1 \right)$, by Chebyshev's inequality we get that $$\mathbb{P}(X \ge \frac{1}{\sqrt{\rho}} \sigma + \mu ) \le \rho$$ 
Unraveling the expression $\frac{1}{\sqrt{\rho}} \sigma + \mu$ with the chosen value of $z$,
we get that it is equal to $r \mu$. Hence, the probability that $X \ge r \mu$ is at most $\rho$, as desired.

Note that the running time of the second act of the algorithm primarily depends only on $z$. Further note that $z = \alpha M/\theta$ where $\alpha$ is a constant depending only on $r, \rho$.
It follows that if we set $\theta = M^{1/2}$, then the second act of the algorithm runs in time $O(M^{1/2})$.

\paragraph{Wrapping up: } Overall, we see that if we set $\theta = M^{1/2}$, then
for any numbers $r, \rho$, our algorithm runs in time $O(d^6 M^{1/2})$ and returns
a value that is at most a factor of $r$ away from the leaves of $\mathcal{D}(P)$ with probability at least $1 - \rho$. Hence, at the cost of a sub-exponential running time,
we obtain a $(r,\rho)$-approximate counter.

% Furthermore, if we are guaranteed that the total number of leaves of $\mathcal{D}(P)$
% is at most a polynomial of the size of $P$, then just setting $\theta$ to be this polynomial value and running just the first 
% act, gives the correct total value
% of $\mathcal{D}(P)$, i.e., we get an estimator with zero variance. The running time in this case is polynomial and is roughly 
% of the order of $O(d^6\mathit{poly}(|P|))$. 

